\begin{TEXT}{Hymna 2025 (Podobenství)}
\SLOKA \Ch{E}{Kaž}dej z nás si hledá svoje \Ch{A}{mís}to svoji \Ch{E}{tvář}\NL
v světě \Ch{c#}{plným} otá\Ch{A}{zek} a odpo\Ch{E}{vědí}\NL
Byl tu jednou jeden kterej \Ch{A}{ne}měl svato\Ch{E}{zář }\NL
ale \Ch{c#}{blíz}ko lidem  \Ch{A}{učil} jak\Ch{H}{ ti} co  \Ch{E}{vědí}
\SLOKA[Ref. A:]  Tak zku\Ch{c#}{s spát} pod hvězdnou \Ch{E}{ oblohou}\NL
a být \Ch{A}{pak} celý den na  \Ch{E}{nohou}\NL
až v nás več\Ch{c#}{er z}mizí vše\Ch{A}{chno} nepřátel\Ch{H}{ství}\NL
pot\Ch{G#}{om za}zn\Ch{c#}{í k }tobě mo\Ch{A}{udrost} po\Ch{H}{dob}en\Ch{E}{ství}\NL
\SLOKA Jak jen úzká cesta vede a bohatství je klam\NL
a základy se mají stavět řádně\NL
a jsou chvíle kdy je třeba dát vše co vůbec mám\NL
a láska platí i pro ty co jsou na dně\NL
\SLOKA[Ref. B:] /: Tak zkus žít pod modrou oblohou\NL
a zvedat ty co už nemohou\NL
až se kolem ohně zrodí společenství\NL
pak snad poznáš sílu co maj podobenství :/ \NL
\end{TEXT}
